




























\begin{document}
\title{Global Well-Posedness for a Semilinear Parabolic Equation with Critical Exponent}
\author{Yang Jianwei}
\date{}

\maketitle
\ifdefined\YTfallback
\noindent\fbox{\parbox{0.96\linewidth}{\small\textbf{[YangTex preview mode] }
The official class file \texttt{amsbook.cls} was not found on this machine, so the
document uses a generic article layout. Put the class file into
\texttt{YangTex/templates/classes/} (or install it with \texttt{tlmgr}),
then recompile to obtain the real journal layout.}}
\bigskip
\fi
\begin{abstract}
We study the Cauchy problem for a semilinear parabolic equation on a bounded domain. By combining energy estimates with a Moser-type iteration we prove the existence of a unique global solution and derive an explicit decay rate.
\end{abstract}
\paragraph{Keywords} semilinear parabolic equation; energy estimate; global existence; Sobolev embedding

\section{Introduction}
Let $\Om\subset\R^n$ be a bounded domain with smooth boundary. We consider the initial-boundary value problem for a semilinear heat equation with a critical growth nonlinearity.

\begin{theorem}
Assume $u_0\in H^1_0(\Om)$ and $1<p<\frac{n+2}{n-2}$. Then there exists a unique global classical solution $u\in C([0,\infty);H^1_0(\Om))$.
\end{theorem}

\begin{proof}
The local existence follows from the contraction mapping principle. A uniform $H^1$ bound obtained from the energy identity excludes finite-time blow-up.
\end{proof}

\section{Main estimates}
\begin{equation}
\frac{\dd}{\dd t}\|u(t)\|_{L^2}^2+2\|\nabla u(t)\|_{L^2}^2 \le C\|u(t)\|_{L^p}^{p},
\end{equation}
The critical case $p=\frac{n+2}{n-2}$ requires a concentration-compactness argument, which will be treated in a forthcoming work.

\section{Concluding remarks}
The method extends to more general domains and boundary conditions; a detailed
treatment will appear elsewhere.

\begin{thebibliography}{9}
\bibitem{Evans} L.~C. Evans, \emph{Partial Differential Equations}, 2nd ed., AMS, 2010.
\bibitem{Brezis} H. Brezis, \emph{Functional Analysis, Sobolev Spaces and PDEs}, Springer, 2011.
\end{thebibliography}
\end{document}
